\documentclass[11pt,a4paper]{article}
\usepackage[T1]{fontenc}
\usepackage{lmodern,amsmath,amssymb,amsthm,mathtools,booktabs,microtype}
\usepackage[margin=26mm]{geometry}
\usepackage[colorlinks=true,linkcolor=blue,citecolor=blue,urlcolor=blue]{hyperref}
\setlength{\emergencystretch}{2em}
\newtheorem{theorem}{Theorem}[section]
\newtheorem{proposition}[theorem]{Proposition}
\newtheorem{lemma}[theorem]{Lemma}
\newtheorem{corollary}[theorem]{Corollary}
\theoremstyle{definition}\newtheorem{definition}[theorem]{Definition}
\theoremstyle{remark}\newtheorem{remark}[theorem]{Remark}
\newcommand{\Q}{\mathbb Q}\newcommand{\Z}{\mathbb Z}
\newcommand{\R}{\mathbb R}\newcommand{\A}{\mathbb A}
\newcommand{\PP}{\mathbb P}\newcommand{\sqcl}[1]{[#1]}
\DeclareMathOperator{\Tr}{Tr}\DeclareMathOperator{\Nm}{N}
\DeclareMathOperator{\N}{N}\DeclareMathOperator{\rad}{rad}
\DeclareMathOperator{\sfpart}{sf}
\title{Four Contacts with Bounded Odd Radical and Unbounded Height}
\author{}\date{16 November 2026}
\begin{document}\maketitle
\begin{abstract}
For the four contacts $\{-1,0,2,3\}$ there are infinitely many
rational arguments whose positive odd contact radical is one outside
$\{2,3,5\}$. Their heights are unbounded, even after restricting
the arguments to $(3,5)$ and excluding the pole at one. An explicit
elliptic curve and a nontorsion rational point prove these assertions
unconditionally. Consequently no uniform logarithmic height bound in
terms of that radical can hold for this contact set.
\end{abstract}

\section{Height and odd contact radicals}
For a reduced fraction $a/b$, with $b>0$, put
\[
                       h(a/b)=\log\max(|a|,b).
\]
For a prime $p$, normalize $v_p$ by $v_p(p)=1$.
Let $\Theta\subset\Q$ be a finite set of distinct contacts, and
let $S$ contain their denominator primes and the primes at which
two contacts have equal reduction. For $t\notin\Theta$, define
\begin{equation}\label{eq:radical}
 M_{\Theta,S}(t)=
 \prod_{\substack{p\notin S:\text{ some }\theta\in\Theta\text{ has}\\
                 v_p(t-\theta)>0\text{ odd}}}p.
\end{equation}
The empty product is one. This note concerns a bound for this radical
alone, with the additional excluded input $t=1$. It does not assert
that the counterexamples satisfy all equations of a Diophantine
construction. Positive bounds with more contacts are given
in~\cite{HeightBounds}.

\section{Four contacts with bounded odd radical}

\begin{theorem}\label{thm:fourcounterexample}
Let
\[
                  \Theta=\{-1,0,2,3\},\qquad S\supseteq\{2,3,5\}
\]
with $S$ finite. There are infinitely many
$t\in\Q\setminus(\Theta\cup\{1\})$ for which
\[
                         M_{\Theta,S}(t)=1.
\]
Their heights are unbounded. Consequently there are no constants
$c,C\ge0$ such that
\[
                h(t)\le c\log M_{\Theta,S}(t)+C
\]
for every rational $t$ outside the contacts and the pole.
No conjecture is used.
\end{theorem}
\begin{proof}
Consider the smooth quartic curve
\begin{equation}\label{eq:quartictwist}
                   10y^2=(t+1)t(t-2)(t-3).
\end{equation}
Its four roots are distinct, and it has the rational point
$(t,y)=(4,2)$. On $t\ne0$, make the substitutions
\[
                        X_E=\frac{60}{t},\qquad
                        Y_E=\frac{600y}{t^2}.
\]
They give the elliptic curve
\begin{equation}\label{eq:longcurve}
 E:\quad Y_E^2=X_E^3+10X_E^2-2400X_E+36000
              =(X_E+60)(X_E-20)(X_E-30),
\end{equation}
with inverse formulas
\[
                         t=\frac{60}{X_E},\qquad
                         y=\frac{6Y_E}{X_E^2}.
\]
The point $(4,2)$ corresponds to $P=(15,75)$.
The tangent slope at $P$ is $-19/2$, and the group law gives
\begin{equation}\label{eq:double}
                         2P=\left(\frac{201}{4},\frac{2079}{8}\right).
\end{equation}

To apply the integral short-form version of Nagell--Lutz, put
\[
                        x=9X_E+30,\qquad Y=27Y_E.
\]
Then
\begin{equation}\label{eq:shortcurve}
                 Y^2=x^3-197100x+32130000.
\end{equation}
This integral equation is nonsingular; its cubic has the distinct
roots $-510,210,300$. The image of $2P$ has abscissa $1929/4$,
which is not integral. Nagell--Lutz states that every rational
torsion point on an integral short Weierstrass equation has
integral coordinates; see~\cite[Theorem~1.1]{Alpoge} or
\cite[Chapter~II, Theorem~5.1]{Milne}. Thus $P$ has infinite order.

For each nonzero integer $n$, the point $[n]P$ is nonidentity
and is not a point of order two. It cannot have $X_E=0$, since
$Y_E^2=36000$ has no rational solution. The inverse formulas
therefore give a rational point $(t_n,y_n)$ on
\eqref{eq:quartictwist} with $t_n\notin\Theta$.
Also $t_n\ne1$, since at $t=1$ the equation would require
$y^2=2/5$. There are infinitely many distinct $t_n$: a fixed
value of $t$ determines at most two points of~\eqref{eq:longcurve},
whereas the multiples of $P$ are all distinct.

Fix $p\notin\{2,3,5\}$. If $v_p(t_n)<0$, then all four
differences $t_n-\theta$ have negative valuation, so $p$
does not occur in~\eqref{eq:radical}. If $v_p(t_n)\ge0$,
at most one difference has positive valuation, since the
nonzero differences among the contacts have prime factors
only two and three. Equation~\eqref{eq:quartictwist} gives
\[
 v_p\bigl((t_n+1)t_n(t_n-2)(t_n-3)\bigr)=2v_p(y_n),
\]
which is even. Therefore no positive contact valuation at such
a prime is odd. It follows that $M_{\Theta,S}(t_n)=1$.

Finally, only finitely many reduced rational fractions have
height at most any fixed constant. Since the $t_n$ form an
infinite set, their heights are unbounded. This proves all
assertions.
\end{proof}

\begin{corollary}\label{cor:positive}
The counterexamples in Theorem~\ref{thm:fourcounterexample}
may be chosen with $3<t<5$, and still have unbounded height.
\end{corollary}
\begin{proof}
The real group $E(\R)$ is compact. The closure of $\Z P$
is an infinite compact subgroup, hence has no isolated point:
otherwise it would be discrete and compact, and therefore finite.
Every sufficiently small neighborhood of $P=(15,75)$ consequently
contains infinitely many multiples of $P$. The rational function
$60/X_E$ is continuous there and equals four at $P$.
Choose the neighborhood so that its values lie in $(3,5)$.
The preceding finite-fiber and height arguments apply to this
infinite subset.
\end{proof}

The theorem concerns a proposed height bound for the odd radical.
It does not assert that these arguments extend to complete
nonintegral witnesses of any Diophantine construction.

\section{Exact certificate}

The standard-library program
\nolinkurl{computation/contact_heights/check_input_height_amplification.py}
checks the rational point and duplication in~\eqref{eq:longcurve},
the birational substitution identity, the integral change to
\eqref{eq:shortcurve}, and the nonintegral doubled abscissa.
It also verifies the distinct roots and the excluded rational
values. The same program checks height constants used in the companion
positive construction~\cite{HeightBounds}.
Its machine-readable result is
\nolinkurl{input-height-amplification-certificate.json} in the
same directory. The infinitude assertion follows from the
proved nontorsion criterion and the argument above; it is not
inferred from a finite search over multiples.

\begin{thebibliography}{99}
\bibitem{Alpoge}
L.~Alpoge, \emph{Nagell--Lutz, quickly}, author manuscript,
Theorem~1.1,
\url{https://people.math.harvard.edu/~alpoge/papers/nagell-lutz,%20quickly.pdf}.

\bibitem{Milne}
J.~S.~Milne, \emph{Elliptic Curves}, second edition,
World Scientific, 2021, Chapter~II, Section~5,
\url{https://www.jmilne.org/math/Books/EC2.pdf}.

\bibitem{HeightBounds}
\emph{Input Height Bounds and Odd Contact Radicals},
project manuscript, 10 November 2026,
\nolinkurl{papers/contact_heights/input_height_amplification_20261110_v22.tex}.
\end{thebibliography}
\end{document}
